\chapter{Data Collapse Method Curves}\label{appendix:data_collapse_figures}

This part of the appendix includes the data collapse curves, as described in section \ref{ssec:data_collapse_method}, for all determined thresholds, 
shown together with the logical error rate curves of the used interval. 

The maximal rejection rate over all distances is shown for the logical error rate curve. 
If not rejection rate curve is shown, then no numerical underflow was detected.

The determined thresholds and the uncertainty on it is shown and plotted. 
Only the value of the determined critical threshold is shown. 

\section{Code-Capacity Noise Model}\label{ssec:fssa_plots_code_capacity}

\begin{figure}
    \centering
    \includegraphics[width=\linewidth]{basic/code_capacity_ml_x_00.pdf}
    \caption{
        This figure illustrates the underlying data and results of the data collapse method for the code-capacity noise model, maximum likelihood decoding and $X$ observable. 
        The left panel shows the data points collapsed onto a single curve using the optimized values of $p_{th}$ and $\nu$, which are also indicated in the plot. 
        The right plot shares the $y$-axis with the left panel. 
        It shows the logical error rate curves over the physical error rate interval used for the data collapse analysis. 
        The determined thresholds together with is uncertainty is marked by the vertical line.
        If numerical underflow occurred, the corresponding rejection rate curve is shown as the faint blue line, referenced to the secondary $y$-axis on the right.  
        }
    \label{fig_fssa:cc_ml_x}
\end{figure}


\begin{figure}
    \centering
    \includegraphics[width=\linewidth]{basic/code_capacity_ml_z_00.pdf}
    \caption{
        This figure illustrates the underlying data and results of the data collapse method for the code-capacity noise model, maximum likelihood decoding and $Z$ observable. 
        The left panel shows the data points collapsed onto a single curve using the optimized values of $p_{th}$ and $\nu$, which are also indicated in the plot. 
        The right plot shares the $y$-axis with the left panel. 
        It shows the logical error rate curves over the physical error rate interval used for the data collapse analysis. 
        The determined thresholds together with is uncertainty is marked by the vertical line.
        If numerical underflow occurred, the corresponding rejection rate curve is shown as the faint blue line, referenced to the secondary $y$-axis on the right.  
        }
    \label{fig_fssa:cc_ml_z}
\end{figure}


\begin{figure}
    \centering
    \includegraphics[width=\linewidth]{basic/code_capacity_ml_z_min_d3_00.pdf}
    \caption{
        This figure illustrates the underlying data and results of the data collapse method for the code-capacity noise model, maximum likelihood decoding and $Z$ observable, with a reduced minimal distance of $d=3$.
        The left panel shows the data points collapsed onto a single curve using the optimized values of $p_{th}$ and $\nu$, which are also indicated in the plot. 
        The right plot shares the $y$-axis with the left panel. 
        It shows the logical error rate curves over the physical error rate interval used for the data collapse analysis. 
        The determined thresholds together with is uncertainty is marked by the vertical line.
        If numerical underflow occurred, the corresponding rejection rate curve is shown as the faint blue line, referenced to the secondary $y$-axis on the right.  
        }
    \label{fig_fssa:cc_ml_z_min_d3}
\end{figure}


\begin{figure}
    \centering
    \includegraphics[width=\linewidth]{basic/code_capacity_mwpm_x_00.pdf}
    \caption{
        This figure illustrates the underlying data and results of the data collapse method for the code-capacity noise model, MWPM decoding and $X$ observable. 
        The left panel shows the data points collapsed onto a single curve using the optimized values of $p_{th}$ and $\nu$, which are also indicated in the plot. 
        The right plot shares the $y$-axis with the left panel. 
        It shows the logical error rate curves over the physical error rate interval used for the data collapse analysis. 
        The determined thresholds together with is uncertainty is marked by the vertical line.
        If numerical underflow occurred, the corresponding rejection rate curve is shown as the faint blue line, referenced to the secondary $y$-axis on the right.  
        }
    \label{fig_fssa:cc_mwpm_x}
\end{figure}


\begin{figure}
    \centering
    \includegraphics[width=\linewidth]{basic/code_capacity_mwpm_z_00.pdf}
    \caption{
        This figure illustrates the underlying data and results of the data collapse method for the code-capacity noise model, MWPM decoding and $Z$ observable. 
        The left panel shows the data points collapsed onto a single curve using the optimized values of $p_{th}$ and $\nu$, which are also indicated in the plot. 
        The right plot shares the $y$-axis with the left panel. 
        It shows the logical error rate curves over the physical error rate interval used for the data collapse analysis. 
        The determined thresholds together with is uncertainty is marked by the vertical line.
        If numerical underflow occurred, the corresponding rejection rate curve is shown as the faint blue line, referenced to the secondary $y$-axis on the right.  
        }
    \label{fig_fssa:cc_mwpm_z}
\end{figure}


% circuit
\clearpage
\section{Circuit-Level Noise Model}
\subsection{Single QEC cycle}\label{ssec:fssa_plots_circ_single_cycle}

\begin{figure}[!b]
    \centering
    \includegraphics[width=\linewidth]{circ/circ_ml_x_00.pdf}
    \caption{
        This figure illustrates the underlying data and results of the data collapse method for the circuit-level noise model, ML decoding and $X$ observable.
        The left panel shows the data points collapsed onto a single curve using the optimized values of $p_{th}$ and $\nu$, which are also indicated in the plot. 
        The right plot shares the $y$-axis with the left panel. 
        It shows the logical error rate curves over the physical error rate interval used for the data collapse analysis. 
        The determined thresholds together with is uncertainty is marked by the vertical line.
        If numerical underflow occurred, the corresponding rejection rate curve is shown as the faint blue line, referenced to the secondary $y$-axis on the right.  
        }
    \label{fig_fssa:circ_ml_x}
\end{figure}

\begin{figure}
    \centering
    \includegraphics[width=\linewidth]{circ/circ_ml_z_00.pdf}
    \caption{
        This figure illustrates the underlying data and results of the data collapse method for the circuit-level noise model, ML decoding and $Z$ observable.
        The left panel shows the data points collapsed onto a single curve using the optimized values of $p_{th}$ and $\nu$, which are also indicated in the plot. 
        The right plot shares the $y$-axis with the left panel. 
        It shows the logical error rate curves over the physical error rate interval used for the data collapse analysis. 
        The determined thresholds together with is uncertainty is marked by the vertical line.
        If numerical underflow occurred, the corresponding rejection rate curve is shown as the faint blue line, referenced to the secondary $y$-axis on the right.  
        }
    \label{fig_fssa:circ_ml_z}
\end{figure}

\begin{figure}
    \centering
    \includegraphics[width=\linewidth]{circ/circ_ml_z_min_d3_00.pdf}
    \caption{
        This figure illustrates the underlying data and results of the data collapse method for the circuit-level noise model, ML decoding and $Z$ observable.
        In contrast to the previous analysis, this data set uses the distance $d=3$ as lower bound for the distance.
        The left panel shows the data points collapsed onto a single curve using the optimized values of $p_{th}$ and $\nu$, which are also indicated in the plot. 
        The right plot shares the $y$-axis with the left panel. 
        It shows the logical error rate curves over the physical error rate interval used for the data collapse analysis. 
        The determined thresholds together with is uncertainty is marked by the vertical line.
        If numerical underflow occurred, the corresponding rejection rate curve is shown as the faint blue line, referenced to the secondary $y$-axis on the right.  
        }
    \label{fig_fssa:circ_ml_z_d3}
\end{figure}

\begin{figure}
    \centering
    \includegraphics[width=\linewidth]{circ/circ_mwpm_x_00.pdf}
    \caption{
        This figure illustrates the underlying data and results of the data collapse method for the circuit-level noise model, MWPM decoding and $X$ observable.
        The left panel shows the data points collapsed onto a single curve using the optimized values of $p_{th}$ and $\nu$, which are also indicated in the plot. 
        The right plot shares the $y$-axis with the left panel. 
        It shows the logical error rate curves over the physical error rate interval used for the data collapse analysis. 
        The determined thresholds together with is uncertainty is marked by the vertical line.
        If numerical underflow occurred, the corresponding rejection rate curve is shown as the faint blue line, referenced to the secondary $y$-axis on the right.  
        }
    \label{fig_fssa:circ_mwpm_x}
\end{figure}

\begin{figure}
    \centering
    \includegraphics[width=\linewidth]{circ/circ_mwpm_z_00.pdf}
    \caption{
        This figure illustrates the underlying data and results of the data collapse method for the circuit-level noise model, MWPM decoding and $Z$ observable.
        The left panel shows the data points collapsed onto a single curve using the optimized values of $p_{th}$ and $\nu$, which are also indicated in the plot. 
        The right plot shares the $y$-axis with the left panel. 
        It shows the logical error rate curves over the physical error rate interval used for the data collapse analysis. 
        The determined thresholds together with is uncertainty is marked by the vertical line.
        If numerical underflow occurred, the corresponding rejection rate curve is shown as the faint blue line, referenced to the secondary $y$-axis on the right.  
        }
    \label{fig_fssa:circ_mwpm_z}
\end{figure}

% circuit ft ml
\begin{figure}
    \centering
    \includegraphics[width=\linewidth]{circ/circ_ml_x_ft_ml_00.pdf}
    \caption{
        This figure illustrates the underlying data and results of the data collapse method for the circuit-level noise model, ML decoding, $X$ observable.
        The residual error was decoded using the efficient ML decoder. 
        The left panel shows the data points collapsed onto a single curve using the optimized values of $p_{th}$ and $\nu$, which are also indicated in the plot. 
        The right plot shares the $y$-axis with the left panel. 
        It shows the logical error rate curves over the physical error rate interval used for the data collapse analysis. 
        The determined thresholds together with is uncertainty is marked by the vertical line.
        If numerical underflow occurred, the corresponding rejection rate curve is shown as the faint blue line, referenced to the secondary $y$-axis on the right.  
        }
    \label{fig_fssa:circ_ml_x_ft}
\end{figure}

\begin{figure}
    \centering
    \includegraphics[width=\linewidth]{circ/circ_ml_z_ft_ml_00.pdf}
    \caption{
        This figure illustrates the underlying data and results of the data collapse method for the circuit-level noise model, ML decoding and $Z$ observable.
        The residual error was decoded using the efficient ML decoder.
        The left panel shows the data points collapsed onto a single curve using the optimized values of $p_{th}$ and $\nu$, which are also indicated in the plot. 
        The right plot shares the $y$-axis with the left panel. 
        It shows the logical error rate curves over the physical error rate interval used for the data collapse analysis. 
        The determined thresholds together with is uncertainty is marked by the vertical line.
        If numerical underflow occurred, the corresponding rejection rate curve is shown as the faint blue line, referenced to the secondary $y$-axis on the right.  
        }
    \label{fig_fssa:circ_ml_z_ft}
\end{figure}

\begin{figure}
    \centering
    \includegraphics[width=\linewidth]{circ/circ_mwpm_x_ft_ml_00.pdf}
    \caption{
        This figure illustrates the underlying data and results of the data collapse method for the circuit-level noise model, MWPM decoding and $X$ observable.
        The residual error was decoded using the efficient ML decoder.
        The left panel shows the data points collapsed onto a single curve using the optimized values of $p_{th}$ and $\nu$, which are also indicated in the plot. 
        The right plot shares the $y$-axis with the left panel. 
        It shows the logical error rate curves over the physical error rate interval used for the data collapse analysis. 
        The determined thresholds together with is uncertainty is marked by the vertical line.
        If numerical underflow occurred, the corresponding rejection rate curve is shown as the faint blue line, referenced to the secondary $y$-axis on the right.  
        }
    \label{fig_fssa:circ_mwpm_x_ft}
\end{figure}

\begin{figure}
    \centering
    \includegraphics[width=\linewidth]{circ/circ_mwpm_z_ft_ml_00.pdf}
    \caption{
        This figure illustrates the underlying data and results of the data collapse method for the circuit-level noise model, MWPM decoding and $Z$ observable.
        The residual error was decoded using the efficient ML decoder.
        The left panel shows the data points collapsed onto a single curve using the optimized values of $p_{th}$ and $\nu$, which are also indicated in the plot. 
        The right plot shares the $y$-axis with the left panel. 
        It shows the logical error rate curves over the physical error rate interval used for the data collapse analysis. 
        The determined thresholds together with is uncertainty is marked by the vertical line.
        If numerical underflow occurred, the corresponding rejection rate curve is shown as the faint blue line, referenced to the secondary $y$-axis on the right.  
        }
    \label{fig_fssa:circ_mwpm_z_ft}
\end{figure}


% Done with ML 
\clearpage
\subsection{Multiple QEC cycles}\label{ssec:fssa_plots_multi_cycle}
% Now multi round

% ML X
\begin{figure}[b]
    \centering
    \includegraphics[width=\linewidth]{multi/multi_round_ml_x_00.pdf}
    \caption{
        This figure illustrates the underlying data and results of the data collapse method for the circuit-level noise model, ML decoding and $X$ observable after $1$ round.
        The left panel shows the data points collapsed onto a single curve using the optimized values of $p_{th}$ and $\nu$, which are also indicated in the plot. 
        The right plot shares the $y$-axis with the left panel. 
        It shows the logical error rate curves over the physical error rate interval used for the data collapse analysis. 
        The determined thresholds together with is uncertainty is marked by the vertical line.
        If numerical underflow occurred, the corresponding rejection rate curve is shown as the faint blue line, referenced to the secondary $y$-axis on the right.  
        }
    \label{fig_fssa:multi_ml_x_r01}
\end{figure}


\begin{figure}
    \centering
    \includegraphics[width=\linewidth]{multi/multi_round_ml_x_01.pdf}
    \caption{
        This figure illustrates the underlying data and results of the data collapse method for the circuit-level noise model, ML decoding and $X$ observable after $2$ rounds.
        The left panel shows the data points collapsed onto a single curve using the optimized values of $p_{th}$ and $\nu$, which are also indicated in the plot. 
        The right plot shares the $y$-axis with the left panel. 
        It shows the logical error rate curves over the physical error rate interval used for the data collapse analysis. 
        The determined thresholds together with is uncertainty is marked by the vertical line.
        If numerical underflow occurred, the corresponding rejection rate curve is shown as the faint blue line, referenced to the secondary $y$-axis on the right.  
        }
    \label{fig_fssa:multi_ml_x_r02}
\end{figure}


\begin{figure}
    \centering
    \includegraphics[width=\linewidth]{multi/multi_round_ml_x_02.pdf}
    \caption{
        This figure illustrates the underlying data and results of the data collapse method for the circuit-level noise model, ML decoding and $X$ observable after $3$ rounds.
        The left panel shows the data points collapsed onto a single curve using the optimized values of $p_{th}$ and $\nu$, which are also indicated in the plot. 
        The right plot shares the $y$-axis with the left panel. 
        It shows the logical error rate curves over the physical error rate interval used for the data collapse analysis. 
        The determined thresholds together with is uncertainty is marked by the vertical line.
        If numerical underflow occurred, the corresponding rejection rate curve is shown as the faint blue line, referenced to the secondary $y$-axis on the right.  
        }
    \label{fig_fssa:multi_ml_x_r03}
\end{figure}


\begin{figure}
    \centering
    \includegraphics[width=\linewidth]{multi/multi_round_ml_x_03.pdf}
    \caption{
        This figure illustrates the underlying data and results of the data collapse method for the circuit-level noise model, ML decoding and $X$ observable after $4$ rounds.
        The left panel shows the data points collapsed onto a single curve using the optimized values of $p_{th}$ and $\nu$, which are also indicated in the plot. 
        The right plot shares the $y$-axis with the left panel. 
        It shows the logical error rate curves over the physical error rate interval used for the data collapse analysis. 
        The determined thresholds together with is uncertainty is marked by the vertical line.
        If numerical underflow occurred, the corresponding rejection rate curve is shown as the faint blue line, referenced to the secondary $y$-axis on the right.  
        }
    \label{fig_fssa:multi_ml_x_r04}
\end{figure}


\begin{figure}
    \centering
    \includegraphics[width=\linewidth]{multi/multi_round_ml_x_04.pdf}
    \caption{
        This figure illustrates the underlying data and results of the data collapse method for the circuit-level noise model, ML decoding and $X$ observable after $6$ rounds.
        The left panel shows the data points collapsed onto a single curve using the optimized values of $p_{th}$ and $\nu$, which are also indicated in the plot. 
        The right plot shares the $y$-axis with the left panel. 
        It shows the logical error rate curves over the physical error rate interval used for the data collapse analysis. 
        The determined thresholds together with is uncertainty is marked by the vertical line.
        If numerical underflow occurred, the corresponding rejection rate curve is shown as the faint blue line, referenced to the secondary $y$-axis on the right.  
        }
    \label{fig_fssa:multi_ml_x_r06}
\end{figure}


\begin{figure}
    \centering
    \includegraphics[width=\linewidth]{multi/multi_round_ml_x_05.pdf}
    \caption{
        This figure illustrates the underlying data and results of the data collapse method for the circuit-level noise model, ML decoding and $X$ observable after $8$ rounds.
        The left panel shows the data points collapsed onto a single curve using the optimized values of $p_{th}$ and $\nu$, which are also indicated in the plot. 
        The right plot shares the $y$-axis with the left panel. 
        It shows the logical error rate curves over the physical error rate interval used for the data collapse analysis. 
        The determined thresholds together with is uncertainty is marked by the vertical line.
        If numerical underflow occurred, the corresponding rejection rate curve is shown as the faint blue line, referenced to the secondary $y$-axis on the right.  
        }
    \label{fig_fssa:multi_ml_x_r08}
\end{figure}


\begin{figure}
    \centering
    \includegraphics[width=\linewidth]{multi/multi_round_ml_x_06.pdf}
    \caption{
        This figure illustrates the underlying data and results of the data collapse method for the circuit-level noise model, ML decoding and $X$ observable after $10$ rounds.
        The left panel shows the data points collapsed onto a single curve using the optimized values of $p_{th}$ and $\nu$, which are also indicated in the plot. 
        The right plot shares the $y$-axis with the left panel. 
        It shows the logical error rate curves over the physical error rate interval used for the data collapse analysis. 
        The determined thresholds together with is uncertainty is marked by the vertical line.
        If numerical underflow occurred, the corresponding rejection rate curve is shown as the faint blue line, referenced to the secondary $y$-axis on the right.  
        }
    \label{fig_fssa:multi_ml_x_r10}
\end{figure}


% ml z


\begin{figure}
    \centering
    \includegraphics[width=\linewidth]{multi/multi_round_ml_z_00.pdf}
    \caption{
        This figure illustrates the underlying data and results of the data collapse method for the circuit-level noise model, ML decoding and $Z$ observable after $1$ rounds.
        The left panel shows the data points collapsed onto a single curve using the optimized values of $p_{th}$ and $\nu$, which are also indicated in the plot. 
        The right plot shares the $y$-axis with the left panel. 
        It shows the logical error rate curves over the physical error rate interval used for the data collapse analysis. 
        The determined thresholds together with is uncertainty is marked by the vertical line.
        If numerical underflow occurred, the corresponding rejection rate curve is shown as the faint blue line, referenced to the secondary $y$-axis on the right.  
        }
    \label{fig_fssa:multi_ml_Z_r01}
\end{figure}


\begin{figure}
    \centering
    \includegraphics[width=\linewidth]{multi/multi_round_ml_z_01.pdf}
    \caption{
        This figure illustrates the underlying data and results of the data collapse method for the circuit-level noise model, ML decoding and $Z$ observable after $2$ rounds.
        The left panel shows the data points collapsed onto a single curve using the optimized values of $p_{th}$ and $\nu$, which are also indicated in the plot. 
        The right plot shares the $y$-axis with the left panel. 
        It shows the logical error rate curves over the physical error rate interval used for the data collapse analysis. 
        The determined thresholds together with is uncertainty is marked by the vertical line.
        If numerical underflow occurred, the corresponding rejection rate curve is shown as the faint blue line, referenced to the secondary $y$-axis on the right.  
        }
    \label{fig_fssa:multi_ml_Z_r02}
\end{figure}


\begin{figure}
    \centering
    \includegraphics[width=\linewidth]{multi/multi_round_ml_z_02.pdf}
    \caption{
        This figure illustrates the underlying data and results of the data collapse method for the circuit-level noise model, ML decoding and $Z$ observable after $3$ rounds.
        The left panel shows the data points collapsed onto a single curve using the optimized values of $p_{th}$ and $\nu$, which are also indicated in the plot. 
        The right plot shares the $y$-axis with the left panel. 
        It shows the logical error rate curves over the physical error rate interval used for the data collapse analysis. 
        The determined thresholds together with is uncertainty is marked by the vertical line.
        If numerical underflow occurred, the corresponding rejection rate curve is shown as the faint blue line, referenced to the secondary $y$-axis on the right.  
        }
    \label{fig_fssa:multi_ml_Z_r03}
\end{figure}


\begin{figure}
    \centering
    \includegraphics[width=\linewidth]{multi/multi_round_ml_z_03.pdf}
    \caption{
        This figure illustrates the underlying data and results of the data collapse method for the circuit-level noise model, ML decoding and $Z$ observable after $4$ rounds.
        The left panel shows the data points collapsed onto a single curve using the optimized values of $p_{th}$ and $\nu$, which are also indicated in the plot. 
        The right plot shares the $y$-axis with the left panel. 
        It shows the logical error rate curves over the physical error rate interval used for the data collapse analysis. 
        The determined thresholds together with is uncertainty is marked by the vertical line.
        If numerical underflow occurred, the corresponding rejection rate curve is shown as the faint blue line, referenced to the secondary $y$-axis on the right.  
        }
    \label{fig_fssa:multi_ml_Z_r04}
\end{figure}


\begin{figure}
    \centering
    \includegraphics[width=\linewidth]{multi/multi_round_ml_z_04.pdf}
    \caption{
        This figure illustrates the underlying data and results of the data collapse method for the circuit-level noise model, ML decoding and $Z$ observable after $6$ rounds.
        The left panel shows the data points collapsed onto a single curve using the optimized values of $p_{th}$ and $\nu$, which are also indicated in the plot. 
        The right plot shares the $y$-axis with the left panel. 
        It shows the logical error rate curves over the physical error rate interval used for the data collapse analysis. 
        The determined thresholds together with is uncertainty is marked by the vertical line.
        If numerical underflow occurred, the corresponding rejection rate curve is shown as the faint blue line, referenced to the secondary $y$-axis on the right.  
        }
    \label{fig_fssa:multi_ml_Z_r06}
\end{figure}


\begin{figure}
    \centering
    \includegraphics[width=\linewidth]{multi/multi_round_ml_z_05.pdf}
    \caption{
        This figure illustrates the underlying data and results of the data collapse method for the circuit-level noise model, ML decoding and $Z$ observable after $8$ rounds.
        The left panel shows the data points collapsed onto a single curve using the optimized values of $p_{th}$ and $\nu$, which are also indicated in the plot. 
        The right plot shares the $y$-axis with the left panel. 
        It shows the logical error rate curves over the physical error rate interval used for the data collapse analysis. 
        The determined thresholds together with is uncertainty is marked by the vertical line.
        If numerical underflow occurred, the corresponding rejection rate curve is shown as the faint blue line, referenced to the secondary $y$-axis on the right.  
        }
    \label{fig_fssa:multi_ml_Z_r08}
\end{figure}


\begin{figure}
    \centering
    \includegraphics[width=\linewidth]{multi/multi_round_ml_z_06.pdf}
    \caption{
        This figure illustrates the underlying data and results of the data collapse method for the circuit-level noise model, ML decoding and $Z$ observable after $10$ rounds.
        The left panel shows the data points collapsed onto a single curve using the optimized values of $p_{th}$ and $\nu$, which are also indicated in the plot. 
        The right plot shares the $y$-axis with the left panel. 
        It shows the logical error rate curves over the physical error rate interval used for the data collapse analysis. 
        The determined thresholds together with is uncertainty is marked by the vertical line.
        If numerical underflow occurred, the corresponding rejection rate curve is shown as the faint blue line, referenced to the secondary $y$-axis on the right.  
        }
    \label{fig_fssa:multi_ml_Z_r10}
\end{figure}

% mwpm x


\begin{figure}
    \centering
    \includegraphics[width=\linewidth]{multi/multi_round_mwpm_x_00.pdf}
    \caption{
        This figure illustrates the underlying data and results of the data collapse method for the circuit-level noise model, MWPM decoding and $X$ observable after $1$ round.
        The left panel shows the data points collapsed onto a single curve using the optimized values of $p_{th}$ and $\nu$, which are also indicated in the plot. 
        The right plot shares the $y$-axis with the left panel. 
        It shows the logical error rate curves over the physical error rate interval used for the data collapse analysis. 
        The determined thresholds together with is uncertainty is marked by the vertical line.
        If numerical underflow occurred, the corresponding rejection rate curve is shown as the faint blue line, referenced to the secondary $y$-axis on the right.  
        }
    \label{fig_fssa:multi_mwpm_x_r01}
\end{figure}


\begin{figure}
    \centering
    \includegraphics[width=\linewidth]{multi/multi_round_mwpm_x_01.pdf}
    \caption{
        This figure illustrates the underlying data and results of the data collapse method for the circuit-level noise model, MWPM decoding and $X$ observable after $2$ rounds.
        The left panel shows the data points collapsed onto a single curve using the optimized values of $p_{th}$ and $\nu$, which are also indicated in the plot. 
        The right plot shares the $y$-axis with the left panel. 
        It shows the logical error rate curves over the physical error rate interval used for the data collapse analysis. 
        The determined thresholds together with is uncertainty is marked by the vertical line.
        If numerical underflow occurred, the corresponding rejection rate curve is shown as the faint blue line, referenced to the secondary $y$-axis on the right.  
        }
    \label{fig_fssa:multi_mwpm_x_r02}
\end{figure}


\begin{figure}
    \centering
    \includegraphics[width=\linewidth]{multi/multi_round_mwpm_x_02.pdf}
    \caption{
        This figure illustrates the underlying data and results of the data collapse method for the circuit-level noise model, MWPM decoding and $X$ observable after $3$ rounds.
        The left panel shows the data points collapsed onto a single curve using the optimized values of $p_{th}$ and $\nu$, which are also indicated in the plot. 
        The right plot shares the $y$-axis with the left panel. 
        It shows the logical error rate curves over the physical error rate interval used for the data collapse analysis. 
        The determined thresholds together with is uncertainty is marked by the vertical line.
        If numerical underflow occurred, the corresponding rejection rate curve is shown as the faint blue line, referenced to the secondary $y$-axis on the right.  
        }
    \label{fig_fssa:multi_mwpm_x_r03}
\end{figure}


\begin{figure}
    \centering
    \includegraphics[width=\linewidth]{multi/multi_round_mwpm_x_03.pdf}
    \caption{
        This figure illustrates the underlying data and results of the data collapse method for the circuit-level noise model, MWPM decoding and $X$ observable after $3$ rounds.
        The left panel shows the data points collapsed onto a single curve using the optimized values of $p_{th}$ and $\nu$, which are also indicated in the plot. 
        The right plot shares the $y$-axis with the left panel. 
        It shows the logical error rate curves over the physical error rate interval used for the data collapse analysis. 
        The determined thresholds together with is uncertainty is marked by the vertical line.
        If numerical underflow occurred, the corresponding rejection rate curve is shown as the faint blue line, referenced to the secondary $y$-axis on the right.  
        }
    \label{fig_fssa:multi_mwpm_x_r04}
\end{figure}


\begin{figure}
    \centering
    \includegraphics[width=\linewidth]{multi/multi_round_mwpm_x_04.pdf}
    \caption{
        This figure illustrates the underlying data and results of the data collapse method for the circuit-level noise model, MWPM decoding and $X$ observable after $4$ rounds.
        The left panel shows the data points collapsed onto a single curve using the optimized values of $p_{th}$ and $\nu$, which are also indicated in the plot. 
        The right plot shares the $y$-axis with the left panel. 
        It shows the logical error rate curves over the physical error rate interval used for the data collapse analysis. 
        The determined thresholds together with is uncertainty is marked by the vertical line.
        If numerical underflow occurred, the corresponding rejection rate curve is shown as the faint blue line, referenced to the secondary $y$-axis on the right.  
        }
    \label{fig_fssa:multi_mwpm_x_r06}
\end{figure}


\begin{figure}
    \centering
    \includegraphics[width=\linewidth]{multi/multi_round_mwpm_x_05.pdf}
    \caption{
        This figure illustrates the underlying data and results of the data collapse method for the circuit-level noise model, MWPM decoding and $X$ observable after $6$ rounds.
        The left panel shows the data points collapsed onto a single curve using the optimized values of $p_{th}$ and $\nu$, which are also indicated in the plot. 
        The right plot shares the $y$-axis with the left panel. 
        It shows the logical error rate curves over the physical error rate interval used for the data collapse analysis. 
        The determined thresholds together with is uncertainty is marked by the vertical line.
        If numerical underflow occurred, the corresponding rejection rate curve is shown as the faint blue line, referenced to the secondary $y$-axis on the right.  
        }
    \label{fig_fssa:multi_mwpm_x_r08}
\end{figure}


\begin{figure}
    \centering
    \includegraphics[width=\linewidth]{multi/multi_round_mwpm_x_06.pdf}
    \caption{
        This figure illustrates the underlying data and results of the data collapse method for the circuit-level noise model, MWPM decoding and $X$ observable after $10$ rounds.
        The left panel shows the data points collapsed onto a single curve using the optimized values of $p_{th}$ and $\nu$, which are also indicated in the plot. 
        The right plot shares the $y$-axis with the left panel. 
        It shows the logical error rate curves over the physical error rate interval used for the data collapse analysis. 
        The determined thresholds together with is uncertainty is marked by the vertical line.
        If numerical underflow occurred, the corresponding rejection rate curve is shown as the faint blue line, referenced to the secondary $y$-axis on the right.  
        }
    \label{fig_fssa:multi_mwpm_x_r10}
\end{figure}

% mwpm z


\begin{figure}
    \centering
    \includegraphics[width=\linewidth]{multi/multi_round_mwpm_z_00.pdf}
    \caption{
        This figure illustrates the underlying data and results of the data collapse method for the circuit-level noise model, MWPM decoding and $X$ observable after $1$ round.
        The left panel shows the data points collapsed onto a single curve using the optimized values of $p_{th}$ and $\nu$, which are also indicated in the plot. 
        The right plot shares the $y$-axis with the left panel. 
        It shows the logical error rate curves over the physical error rate interval used for the data collapse analysis. 
        The determined thresholds together with is uncertainty is marked by the vertical line.
        If numerical underflow occurred, the corresponding rejection rate curve is shown as the faint blue line, referenced to the secondary $y$-axis on the right.  
        }
    \label{fig_fssa:multi_mwpm_z_r01}
\end{figure}


\begin{figure}
    \centering
    \includegraphics[width=\linewidth]{multi/multi_round_mwpm_x_01.pdf}
    \caption{
        This figure illustrates the underlying data and results of the data collapse method for the circuit-level noise model, MWPM decoding and $X$ observable after $2$ rounds.
        The left panel shows the data points collapsed onto a single curve using the optimized values of $p_{th}$ and $\nu$, which are also indicated in the plot. 
        The right plot shares the $y$-axis with the left panel. 
        It shows the logical error rate curves over the physical error rate interval used for the data collapse analysis. 
        The determined thresholds together with is uncertainty is marked by the vertical line.
        If numerical underflow occurred, the corresponding rejection rate curve is shown as the faint blue line, referenced to the secondary $y$-axis on the right.  
        }
    \label{fig_fssa:multi_mwpm_z_r02}
\end{figure}


\begin{figure}
    \centering
    \includegraphics[width=\linewidth]{multi/multi_round_mwpm_x_02.pdf}
    \caption{
        This figure illustrates the underlying data and results of the data collapse method for the circuit-level noise model, MWPM decoding and $X$ observable after $3$ rounds.
        The left panel shows the data points collapsed onto a single curve using the optimized values of $p_{th}$ and $\nu$, which are also indicated in the plot. 
        The right plot shares the $y$-axis with the left panel. 
        It shows the logical error rate curves over the physical error rate interval used for the data collapse analysis. 
        The determined thresholds together with is uncertainty is marked by the vertical line.
        If numerical underflow occurred, the corresponding rejection rate curve is shown as the faint blue line, referenced to the secondary $y$-axis on the right.  
        }
    \label{fig_fssa:multi_mwpm_z_r03}
\end{figure}


\begin{figure}
    \centering
    \includegraphics[width=\linewidth]{multi/multi_round_mwpm_x_03.pdf}
    \caption{
        This figure illustrates the underlying data and results of the data collapse method for the circuit-level noise model, MWPM decoding and $X$ observable after $4$ rounds.
        The left panel shows the data points collapsed onto a single curve using the optimized values of $p_{th}$ and $\nu$, which are also indicated in the plot. 
        The right plot shares the $y$-axis with the left panel. 
        It shows the logical error rate curves over the physical error rate interval used for the data collapse analysis. 
        The determined thresholds together with is uncertainty is marked by the vertical line.
        If numerical underflow occurred, the corresponding rejection rate curve is shown as the faint blue line, referenced to the secondary $y$-axis on the right.  
        }
    \label{fig_fssa:multi_mwpm_z_r04}
\end{figure}


\begin{figure}
    \centering
    \includegraphics[width=\linewidth]{multi/multi_round_mwpm_x_04.pdf}
    \caption{
        This figure illustrates the underlying data and results of the data collapse method for the circuit-level noise model, MWPM decoding and $X$ observable after $6$ rounds.
        The left panel shows the data points collapsed onto a single curve using the optimized values of $p_{th}$ and $\nu$, which are also indicated in the plot. 
        The right plot shares the $y$-axis with the left panel. 
        It shows the logical error rate curves over the physical error rate interval used for the data collapse analysis. 
        The determined thresholds together with is uncertainty is marked by the vertical line.
        If numerical underflow occurred, the corresponding rejection rate curve is shown as the faint blue line, referenced to the secondary $y$-axis on the right.  
        }
    \label{fig_fssa:multi_mwpm_z_r06}
\end{figure}


\begin{figure}
    \centering
    \includegraphics[width=\linewidth]{multi/multi_round_mwpm_x_05.pdf}
    \caption{
        This figure illustrates the underlying data and results of the data collapse method for the circuit-level noise model, MWPM decoding and $X$ observable after $8$ rounds.
        The left panel shows the data points collapsed onto a single curve using the optimized values of $p_{th}$ and $\nu$, which are also indicated in the plot. 
        The right plot shares the $y$-axis with the left panel. 
        It shows the logical error rate curves over the physical error rate interval used for the data collapse analysis. 
        The determined thresholds together with is uncertainty is marked by the vertical line.
        If numerical underflow occurred, the corresponding rejection rate curve is shown as the faint blue line, referenced to the secondary $y$-axis on the right.  
        }
    \label{fig_fssa:multi_mwpm_z_r08}
\end{figure}


\begin{figure}
    \centering
    \includegraphics[width=\linewidth]{multi/multi_round_mwpm_x_06.pdf}
    \caption{
        This figure illustrates the underlying data and results of the data collapse method for the circuit-level noise model, MWPM decoding and $X$ observable after $10$ rounds.
        The left panel shows the data points collapsed onto a single curve using the optimized values of $p_{th}$ and $\nu$, which are also indicated in the plot. 
        The right plot shares the $y$-axis with the left panel. 
        It shows the logical error rate curves over the physical error rate interval used for the data collapse analysis. 
        The determined thresholds together with is uncertainty is marked by the vertical line.
        If numerical underflow occurred, the corresponding rejection rate curve is shown as the faint blue line, referenced to the secondary $y$-axis on the right.  
        }
    \label{fig_fssa:multi_mwpm_z_r10}
\end{figure}

% reference X
\clearpage
\subsection{Conventional Syndrome Extraction Circuit}\label{ssec:fssa_plots_cse}

\begin{figure}
    \centering
    \includegraphics[width=\linewidth]{multi/multi_round_mwpm_x_00.pdf}
    \caption{
        This figure illustrates the underlying data and results of the data collapse method for the circuit-level noise model 
        using the conventional syndrome extraction circuit, MWPM decoding and $X$ observable after $1$ round.
        The left panel shows the data points collapsed onto a single curve using the optimized values of $p_{th}$ and $\nu$, which are also indicated in the plot. 
        The right plot shares the $y$-axis with the left panel. 
        It shows the logical error rate curves over the physical error rate interval used for the data collapse analysis. 
        The determined thresholds together with is uncertainty is marked by the vertical line.
        If numerical underflow occurred, the corresponding rejection rate curve is shown as the faint blue line, referenced to the secondary $y$-axis on the right.  
        }
    \label{fig_fssa:multi_ref_x_r01}
\end{figure}

\begin{figure}
    \centering
    \includegraphics[width=\linewidth]{multi/multi_round_mwpm_x_01.pdf}
    \caption{
        This figure illustrates the underlying data and results of the data collapse method for the circuit-level noise model 
        using the conventional syndrome extraction circuit, MWPM decoding and $X$ observable after $2$ rounds.
        The left panel shows the data points collapsed onto a single curve using the optimized values of $p_{th}$ and $\nu$, which are also indicated in the plot. 
        The right plot shares the $y$-axis with the left panel. 
        It shows the logical error rate curves over the physical error rate interval used for the data collapse analysis. 
        The determined thresholds together with is uncertainty is marked by the vertical line.
        If numerical underflow occurred, the corresponding rejection rate curve is shown as the faint blue line, referenced to the secondary $y$-axis on the right.  
        }
    \label{fig_fssa:multi_ref_x_r02}
\end{figure}

\begin{figure}
    \centering
    \includegraphics[width=\linewidth]{multi/multi_round_mwpm_x_02.pdf}
    \caption{
        This figure illustrates the underlying data and results of the data collapse method for the circuit-level noise model 
        using the conventional syndrome extraction circuit, MWPM decoding and $X$ observable after $4$ rounds.
        The left panel shows the data points collapsed onto a single curve using the optimized values of $p_{th}$ and $\nu$, which are also indicated in the plot. 
        The right plot shares the $y$-axis with the left panel. 
        It shows the logical error rate curves over the physical error rate interval used for the data collapse analysis. 
        The determined thresholds together with is uncertainty is marked by the vertical line.
        If numerical underflow occurred, the corresponding rejection rate curve is shown as the faint blue line, referenced to the secondary $y$-axis on the right.  
        }
    \label{fig_fssa:multi_ref_x_r04}
\end{figure}

\begin{figure}
    \centering
    \includegraphics[width=\linewidth]{multi/multi_round_mwpm_x_03.pdf}
    \caption{
        This figure illustrates the underlying data and results of the data collapse method for the circuit-level noise model 
        using the conventional syndrome extraction circuit, MWPM decoding and $X$ observable after $6$ rounds.
        The left panel shows the data points collapsed onto a single curve using the optimized values of $p_{th}$ and $\nu$, which are also indicated in the plot. 
        The right plot shares the $y$-axis with the left panel. 
        It shows the logical error rate curves over the physical error rate interval used for the data collapse analysis. 
        The determined thresholds together with is uncertainty is marked by the vertical line.
        If numerical underflow occurred, the corresponding rejection rate curve is shown as the faint blue line, referenced to the secondary $y$-axis on the right.  
        }
    \label{fig_fssa:multi_ref_x_r06}
\end{figure}
\clearpage
% Z
\begin{figure}
    \centering
    \includegraphics[width=\linewidth]{multi/multi_round_mwpm_z_00.pdf}
    \caption{
        This figure illustrates the underlying data and results of the data collapse method for the circuit-level noise model 
        using the conventional syndrome extraction circuit, MWPM decoding and $Z$ observable after $1$ round.
        The left panel shows the data points collapsed onto a single curve using the optimized values of $p_{th}$ and $\nu$, which are also indicated in the plot. 
        The right plot shares the $y$-axis with the left panel. 
        It shows the logical error rate curves over the physical error rate interval used for the data collapse analysis. 
        The determined thresholds together with is uncertainty is marked by the vertical line.
        If numerical underflow occurred, the corresponding rejection rate curve is shown as the faint blue line, referenced to the secondary $y$-axis on the right.  
        }
    \label{fig_fssa:multi_ref_z_r01}
\end{figure}
\begin{figure}
    \centering
    \includegraphics[width=\linewidth]{multi/multi_round_mwpm_z_01.pdf}
    \caption{
        This figure illustrates the underlying data and results of the data collapse method for the circuit-level noise model 
        using the conventional syndrome extraction circuit, MWPM decoding and $Z$ observable after $2$ rounds.
        The left panel shows the data points collapsed onto a single curve using the optimized values of $p_{th}$ and $\nu$, which are also indicated in the plot. 
        The right plot shares the $y$-axis with the left panel. 
        It shows the logical error rate curves over the physical error rate interval used for the data collapse analysis. 
        The determined thresholds together with is uncertainty is marked by the vertical line.
        If numerical underflow occurred, the corresponding rejection rate curve is shown as the faint blue line, referenced to the secondary $y$-axis on the right.  
        }
    \label{fig_fssa:multi_ref_z_r02}
\end{figure}

\begin{figure}
    \centering
    \includegraphics[width=\linewidth]{multi/multi_round_mwpm_z_02.pdf}
    \caption{
        This figure illustrates the underlying data and results of the data collapse method for the circuit-level noise model 
        using the conventional syndrome extraction circuit, MWPM decoding and $Z$ observable after $4$ rounds.
        The left panel shows the data points collapsed onto a single curve using the optimized values of $p_{th}$ and $\nu$, which are also indicated in the plot. 
        The right plot shares the $y$-axis with the left panel. 
        It shows the logical error rate curves over the physical error rate interval used for the data collapse analysis. 
        The determined thresholds together with is uncertainty is marked by the vertical line.
        If numerical underflow occurred, the corresponding rejection rate curve is shown as the faint blue line, referenced to the secondary $y$-axis on the right.  
        }
    \label{fig_fssa:multi_ref_z_r04}
\end{figure}
\begin{figure}
    \centering
    \includegraphics[width=\linewidth]{multi/multi_round_mwpm_z_03.pdf}
    \caption{
        This figure illustrates the underlying data and results of the data collapse method for the circuit-level noise model 
        using the conventional syndrome extraction circuit, MWPM decoding and $Z$ observable after $6$ rounds.
        The left panel shows the data points collapsed onto a single curve using the optimized values of $p_{th}$ and $\nu$, which are also indicated in the plot. 
        The right plot shares the $y$-axis with the left panel. 
        It shows the logical error rate curves over the physical error rate interval used for the data collapse analysis. 
        The determined thresholds together with is uncertainty is marked by the vertical line.
        If numerical underflow occurred, the corresponding rejection rate curve is shown as the faint blue line, referenced to the secondary $y$-axis on the right.  
        }
    \label{fig_fssa:multi_ref_z_r06}
\end{figure}